\documentclass[11pt,a4paper]{article}
\usepackage[utf8]{inputenc}
\usepackage[vietnamese]{babel}
\usepackage{amsmath,amssymb}
\usepackage{geometry}
\geometry{a4paper, total={180mm,260mm}, left=15mm, right=15mm, top=18mm, bottom=18mm}
\usepackage{tikz}
\usetikzlibrary{shapes.geometric, arrows.meta, positioning, calc}
\usepackage{tcolorbox}
\tcbuselibrary{skins,breakable}
\usepackage{xcolor}
\usepackage{multicol}
\usepackage{array}
\usepackage{fontawesome}
\usepackage{fancyhdr}

\pagestyle{fancy}
\fancyhf{}
\renewcommand{\headrulewidth}{0.5pt}
\renewcommand{\footrulewidth}{0.5pt}
\fancyhead[L]{\small\bfseries\color{teal}\faBook\ SỔ TAY KIẾN THỨC TOÁN 3}
\fancyhead[R]{\small\bfseries\color{teal}Lớp Toán Thầy Lĩnh -- loptoanthaylinh.vn}
\fancyfoot[C]{\small\color{gray}-- Sổ tay kiến thức: Bài 50. Chu vi hình tam giác, tứ giác, hình chữ nhật, hình vuông --}

\definecolor{primary}{HTML}{005F73}
\definecolor{secondary}{HTML}{0A9396}
\definecolor{accent}{HTML}{EE9B00}
\definecolor{danger}{HTML}{AE2012}
\definecolor{boxbg}{HTML}{F4F6F8}

\begin{document}

% --- BANNER TIÊU ĐỀ ---
\begin{tcolorbox}[
  enhanced,
  colback=white,
  colframe=primary,
  arc=3mm,
  boxrule=1.2pt,
  left=4mm, right=4mm, top=2.5mm, bottom=2.5mm,
  drop shadow=black!10!white
]
  \centering
  {\Large\bfseries\color{primary}\faBookmark\ BÀI 50: CHU VI HÌNH TAM GIÁC, TỨ GIÁC, CHỮ NHẬT, VUÔNG}\\[1.5mm]
  {\small\itshape\color{black!75}Khái niệm chu vi là tổng độ dài các cạnh bao quanh, công thức tính chu vi và bài toán hàng rào trên 1 trang A4}
\end{tcolorbox}

\vspace{-2mm}


\noindent
\begin{minipage}[t]{0.49\textwidth}\vspace{0pt}
\begin{tcolorbox}[
  enhanced,
  colback=blue!4!white,
  colframe=blue!75!black,
  arc=2mm, boxrule=0.8pt,
  left=2mm, right=2mm, top=1.8mm, bottom=1.8mm,
  title={\footnotesize\bfseries\color{white}\faKey\ 1. Chu vi tam giác và tứ giác}
]
{\footnotesize
\begin{itemize}\setlength{\itemsep}{1.5pt}\setlength{\leftskip}{-2.5mm}
  \item \textbf{Khái niệm chu vi:} Chu vi của một hình là \textbf{tổng độ dài} các cạnh bao quanh hình đó.
  \item \textbf{Chu vi hình tam giác:}
  $$P = a + b + c$$
  (Nếu tam giác có 3 cạnh bằng nhau: $P = a \times 3$).
  \item \textbf{Chu vi hình tứ giác:}
  $$P = a + b + c + d$$
  \item Bắt buộc các cạnh phải cùng một \textbf{đơn vị đo}!
\end{itemize}
}
\end{tcolorbox}
\end{minipage}
\hfill
\begin{minipage}[t]{0.49\textwidth}\vspace{0pt}
\begin{tcolorbox}[
  enhanced,
  colback=teal!4!white,
  colframe=teal!75!black,
  arc=2mm, boxrule=0.8pt,
  left=2mm, right=2mm, top=1.8mm, bottom=1.8mm,
  title={\footnotesize\bfseries\color{white}\faCheckSquare\ 2. Chu vi hình chữ nhật và hình vuông}
]
{\footnotesize
\begin{itemize}\setlength{\itemsep}{1.5pt}\setlength{\leftskip}{-2.5mm}
  \item \textbf{Chu vi hình chữ nhật ($a$ là dài, $b$ là rộng):}
  $$\mathbf{P = (a + b) \times 2}$$
  $$\text{Nửa chu vi} = a + b = P : 2$$
  \item \textbf{Chu vi hình vuông ($a$ là độ dài cạnh):}
  $$\mathbf{P = a \times 4} \iff a = P : 4$$
  \item Muốn tìm cạnh hình vuông khi biết chu vi, ta lấy \textbf{chu vi chia cho 4}.
\end{itemize}
}
\end{tcolorbox}
\end{minipage}

\vspace{-2.5mm}
% --- KHỐI 3: SƠ ĐỒ TƯ DUY TIKZ ---
\noindent
\begin{tcolorbox}[
  enhanced,
  colback=orange!4!white,
  colframe=orange!80!black,
  arc=2mm, boxrule=0.8pt,
  left=2mm, right=2mm, top=1.5mm, bottom=1.5mm,
  title={\footnotesize\bfseries\color{white}\faSitemap\ SƠ ĐỒ TƯ DUY: HỆ THỐNG CÔNG THỨC CHU VI 4 HÌNH CƠ BẢN}
]
\centering
\begin{tikzpicture}[
  box/.style={draw, rounded corners=2pt, align=center, font=\scriptsize\bfseries, inner sep=2.5pt, minimum height=0.75cm},
  arr/.style={-{Stealth[length=2mm]}, thick}
]
  \node[box, fill=purple!15, draw=purple!80!black, text width=2.6cm] (p1) at (0, 0) {TAM GIÁC\\$P = a + b + c$\\Đều: $P = a \times 3$};
  \node[box, fill=teal!15, draw=teal!80!black, text width=2.6cm] (p2) at (3.5, 0) {TỨ GIÁC\\$P = a + b + c + d$\\Tổng 4 cạnh};
  \node[box, fill=blue!12, draw=blue!80!black, text width=2.6cm] (p3) at (7.0, 0) {HÌNH CHỮ NHẬT\\$\mathbf{P = (a + b) \times 2}$\\Nửa chu vi $= P : 2$};
  \node[box, fill=orange!15, draw=orange!80!black, text width=2.6cm] (p4) at (10.5, 0) {HÌNH VUÔNG\\$\mathbf{P = a \times 4}$\\Cạnh $a = P : 4$};

  \draw[arr, purple!80!black] (p1) -- (p2);
  \draw[arr, teal!80!black] (p2) -- (p3);
  \draw[arr, blue!80!black] (p3) -- (p4);
\end{tikzpicture}
\end{tcolorbox}

\vspace{-2.5mm}
% --- KHỐI 4: VÍ DỤ & MẸO TRÁNH BẪY ---
\noindent
\begin{minipage}[t]{0.49\textwidth}\vspace{0pt}
\begin{tcolorbox}[
  enhanced,
  colback=green!4!white,
  colframe=green!60!black,
  arc=2mm, boxrule=0.8pt,
  left=2mm, right=2mm, top=1.8mm, bottom=1.8mm,
  title={\footnotesize\bfseries\color{white}\faCheckCircle\ Ví dụ thực hành chuẩn mực}
]
{\footnotesize
\textbf{1. Chu vi hình chữ nhật:}\\
Mảnh vườn hình chữ nhật có chiều dài $25\text{ m}$, chiều rộng $15\text{ m}$.\\
Chu vi mảnh vườn: $(25 + 15) \times 2 = \mathbf{80\text{ m}}$.\\
\textbf{2. Tìm cạnh hình vuông:}\\
Một khung tranh hình vuông có chu vi $64\text{ cm}$.\\
Độ dài cạnh khung tranh là: $64 : 4 = \mathbf{16\text{ cm}}$.\\
\textbf{3. Bài toán trừ cổng hàng rào:}\\
Rào khu đất chu vi $80\text{ m}$, để cổng rộng $3\text{ m}$. Chiều dài hàng rào là: $80 - 3 = \mathbf{77\text{ m}}$.
}
\end{tcolorbox}
\end{minipage}
\hfill
\begin{minipage}[t]{0.49\textwidth}\vspace{0pt}
\begin{tcolorbox}[
  enhanced,
  colback=red!4!white,
  colframe=red!70!black,
  arc=2mm, boxrule=0.8pt,
  left=2mm, right=2mm, top=1.8mm, bottom=1.8mm,
  title={\footnotesize\bfseries\color{white}\faExclamationTriangle\ Mẹo tránh bẫy trắc nghiệm}
]
{\footnotesize
\begin{itemize}\setlength{\itemsep}{1.5pt}\setlength{\leftskip}{-2.5mm}
  \item \textbf{Bẫy quên nhân 2 hoặc quên đóng mở ngoặc:}\\
  Công thức tính chu vi hình chữ nhật là $(a + b) \times 2$. Rất nhiều bạn bấm máy hoặc viết $a + b \times 2$ dẫn đến kết quả sai!
  \item \textbf{Bẫy chưa cùng đơn vị đo:}\\
  Chiều dài $2\text{ dm}$, chiều rộng $15\text{ cm}$ $\implies$ phải đổi $2\text{ dm} = 20\text{ cm}$ trước khi tính chu vi!
\end{itemize}
}
\end{tcolorbox}
\end{minipage}


\end{document}
